\documentclass[11pt,a4paper]{article}
\usepackage[T1]{fontenc}
\usepackage[utf8]{inputenc}
\usepackage{lmodern}
\usepackage[margin=2.5cm]{geometry}
\usepackage{amsmath,amssymb}
\usepackage{booktabs}
\usepackage{graphicx}
\usepackage[hidelinks]{hyperref}
\hypersetup{pdfauthor={Gianni Pilotti},pdftitle={Liquidity and Price Impact in Two Uniswap Pools},pdfcreator={},pdfproducer={}}
\setlength{\parindent}{0pt}
\setlength{\parskip}{6pt}

\title{Liquidity and Price Impact in Two Uniswap Pools}
\author{Gianni Pilotti}
\date{}

\begin{document}
\maketitle

\section{The Core Idea}

Market liquidity is the cost of trading a given size. On an automated market maker (AMM) the cost is set by the pool reserves, so it can be read from the chain. The project compares two pools that sit at opposite ends of the scale: a new, almost empty pool (HLD/ETH, Uniswap V4 on Base) and a deep one (ETH/USDC, Uniswap V3 on Ethereum).

The two pools need two different treatments:

\begin{itemize}
\item \textbf{HLD/ETH is modelled:} it has almost no trades, so returns and volatility cannot be estimated. Its costs are computed from the real reserves with the constant-product formula.
\item \textbf{ETH/USDC is estimated:} it has thousands of trades, so price impact is measured from real swap events.
\end{itemize}

\section{Method}

\subsection{Reading the V4 Pool State}

A Uniswap V4 pool is identified by the keccak256 hash of its key (two currencies, fee, tick spacing, hooks). The pool state sits in a mapping of the \texttt{PoolManager} contract, so the code computes the storage slot and reads it with \texttt{extsload}. The first word packs four fields: the square-root price $\sqrt P\cdot2^{96}$ in the low 160 bits, then the tick (24 bits), the protocol fee (24 bits) and the LP fee (24 bits). The code decodes them with bit masks. The liquidity $L$ is read from a second slot.

\subsection{Constant-Product Model}

From $L$ and $\sqrt P$ the virtual reserves of the pool are
\begin{equation}
x=\frac{L}{\sqrt P},\qquad y=L\sqrt P,\qquad k=xy,
\end{equation}
with $x$ in ETH and $y$ in HLD. For an order of $\Delta x$ ETH and a fee $\phi$, the code (\texttt{src/amm.py}) computes
\begin{equation}
x'=x+\Delta x(1-\phi),\qquad y'=\frac{k}{x'},\qquad \Delta y=y-y',
\end{equation}
\begin{equation}
\text{slippage}=\frac{\Delta x/\Delta y}{x/y}-1 .
\end{equation}
The ETH needed to push the price up by a fraction $p$ is $x(\sqrt{1+p}-1)/(1-\phi)$. The impermanent loss of a 50/50 position, for a price ratio $r$, is $2\sqrt r/(1+r)-1$.

\subsection{Estimators on Real Swaps}

For the ETH/USDC swaps (\texttt{src/microstructure.py}), with $q_i$ the signed ETH flow of swap $i$ (positive when the taker buys ETH):
\begin{itemize}
\item \textbf{Kyle's lambda:} the OLS slope of the price change on the signed flow, $\Delta p_i=\alpha+\lambda q_i+\varepsilon_i$, between consecutive swaps.
\item \textbf{Realised volatility:} the standard deviation of log price changes per swap, scaled to a year with the average time between swaps $\bar{dt}$:
\end{itemize}
\begin{equation}
\hat\sigma=s_{\Delta\ln p}\sqrt{\frac{365\cdot24\cdot3600}{\bar{dt}}}.
\end{equation}
\begin{itemize}
\item \textbf{Effective spread:} for each swap, twice the gap between the execution price and the mid price before the swap, in basis points:
\end{itemize}
\begin{equation}
\text{ES}_i=2\,\frac{|P^{\text{exec}}_i-P^{\text{mid}}_{i-1}|}{P^{\text{mid}}_{i-1}}\times10^4 .
\end{equation}
The pool quotes WETH per USDC, so the fitted $\lambda$ is converted to dollars per ETH with $\lambda_{\$}=|\lambda|\,P^2$, where $P$ is the ETH price in dollars (since $P=1/p$ and $dP=-P^2dp$).

\section{Results}

HLD/ETH uses the pool state at Base block 49,023,454. ETH/USDC uses 6,737 swaps over 36 hours (22--23 July 2026), with an average of 28.5 seconds between swaps.

\begin{center}
\begin{tabular}{lll}
\toprule
Metric & HLD/ETH (modelled) & ETH/USDC (estimated) \\
\midrule
Trades observed & about 0 & 6,737 in 36 h \\
Cost of a small order & 1 EUR order: 2,126 bp & 10.2 bp (median effective spread) \\
Capital to move the price $+1\%$ & $1.5\times10^{-5}$ ETH (about 0.025 EUR) & 1,618 ETH (about \$3.05M) \\
Kyle's lambda & not estimable & 0.0116 \$/ETH per ETH \\
$t$-statistic, $R^2$ of the fit & not applicable & 317, 0.94 \\
Realised volatility & not identifiable & 24.2\% annualised \\
\bottomrule
\end{tabular}
\end{center}

\begin{itemize}
\item \textbf{HLD/ETH:} the pool value is about 10 EUR (twice the 0.003 ETH reserve at 1,645.88 EUR per ETH). The LP fee is 1\%.
\item \textbf{ETH/USDC:} the mean effective spread is 15.7 bp, above the median of 10.2 bp, so the distribution is skewed to the right.
\item \textbf{Depth gap:} the capital needed to move the price by 1\% differs by a factor of about $10^8$ between the pools.
\end{itemize}

\section{Worked Example}

\textbf{A 1 EUR order in HLD/ETH.} The reserves are $x=0.003$ ETH and $y=65{,}992{,}981$ HLD, so $k=197{,}978.94$.

\textbf{Step 1.} $\Delta x=1/1645.88=6.0758\times10^{-4}$ ETH.

\textbf{Step 2.} After the 1\% fee, $0.99\cdot6.0758\times10^{-4}=6.0150\times10^{-4}$ ETH enters the curve, so $x'=0.0036015$ and $y'=197{,}978.94/0.0036015=54{,}971{,}217$.

\textbf{Step 3.} The buyer receives $\Delta y=65{,}992{,}981-54{,}971{,}217=11{,}021{,}764$ HLD.

\textbf{Step 4.} The execution price is $6.0758\times10^{-4}/11{,}021{,}764=5.5125\times10^{-11}$ ETH per HLD against a spot of $4.5459\times10^{-11}$. The slippage is $+21.26\%$, or 2,126 bp. The order takes 16.7\% of the HLD reserve, which explains the size of the cost.

\textbf{Moving the price by 1\%.} $x(\sqrt{1.01}-1)/0.99=1.511\times10^{-5}$ ETH, which is about 0.025 EUR at 1,645.88 EUR per ETH.

\textbf{The same move in ETH/USDC.} With $|\lambda|=3.2834\times10^{-9}$ WETH per USDC per ETH, $\lambda_{\$}=3.2834\times10^{-9}\cdot1881.94^2=0.011629$ \$/ETH per ETH. A 1\% move is \$18.82, which needs $18.82/0.011629=1{,}618$ ETH, or \$3.05M. A trade of 10 ETH moves the price by about \$0.116.

\section{Limits}

\begin{itemize}
\item \textbf{Local approximation:} constant product is only a local description of V4 concentrated liquidity. Far from the current price the real curve differs.
\item \textbf{Modelled costs:} the HLD/ETH figures come from a formula and one block of data, not from observed trades. Reserves are virtual, not contract balances.
\item \textbf{Linear extrapolation:} the \$3.05M figure assumes the impact stays linear for a move far larger than the observed 5-minute moves. It gives an order of magnitude only.
\item \textbf{Short sample:} 36 hours of ETH/USDC data say nothing about other market regimes.
\item \textbf{Measurement:} the effective spread uses the price before the swap as mid, so it mixes the pool fee, the price impact of the swap itself and any price change between swaps.
\end{itemize}

\end{document}
